% Generated from locale/tr/content/first-order-logic/syntax-and-semantics/structures.tex; do not hand-edit.


\olfileid[tr]{fol}{syn}{str}

\olsection{Birinci Dereceden Diller İçin \printtoken{P}{structure}}

\begin{explain}
Birinci dereceden diller kendi başlarına \emph{yorumlanmamıştır}:
!!{constant}s, !!{function}s ve !!{predicate}s ifadelerine belirli bir
anlam iliştirilmez. Anlamlar, bir \article{structure}
\emph{!!{structure}} belirtilerek verilir. Bu yapı \emph{evreni}, yani
!!{constant}s ifadelerinin gösterdiği, !!{function}s ifadelerinin
üzerinde işlem yaptığı ve niceleyicilerin üzerinde dolaştığı nesneleri
belirler. Ayrıca hangi !!{constant}s ifadelerinin hangi nesneleri
gösterdiğini, !!a{function} ifadesinin nesneleri nesnelere nasıl
eşlediğini ve !!{predicate}s ifadelerinin hangi nesnelere uygulandığını
belirler. !!^{structure}s, mantıktaki \emph{semantik} kavramların;
örneğin sonuç, geçerlilik ve sağlanabilirlik kavramlarının temelidir.
Literatürde bunlara farklı biçimlerde ``yapılar'', ``yorumlar'' ya da
``modeller'' denir.
\end{explain}

\begin{defn}[!!^{structure}s]
Bir \Article{structure} \emph{!!{structure}}~$\Struct M$, birinci
dereceden mantığın~$\Lang L$ dili için aşağıdaki öğelerden oluşur:
\begin{enumerate}
\item \emph{Evren:} boş olmayan bir küme, $\Domain M$.
\item \emph{!!{constant}s yorumu:} $\Lang L$ dilindeki her
  !!{constant}~$c$ için !!a{element}~$\Assign cM\in\Domain M$.
\item \emph{!!{predicate}s yorumu:} $\Lang L$ dilindeki ($\eq$
  dışındaki) her $n$-li !!{predicate}~$R$ için $n$-li bir bağıntı
  $\Assign RM\subseteq\Domain M^n$.
\item \emph{!!{function}s yorumu:} $\Lang L$ dilindeki her $n$-li
  !!{function}~$f$ için $n$-li bir fonksiyon
  $\Assign fM\colon\Domain M^n\to\Domain M$.
\end{enumerate}
\end{defn}

\begin{ex}
Aritmetik diline ilişkin !!^a{structure}~$\Struct M$ şunlardan oluşur:
bir küme; !!{constant}~$\Obj0$ simgesinin yorumu olarak
$\Domain M$ içindeki $\Assign{\Obj0}M$ elemanı; bir-li
$\Assign{\Obj\prime}M\colon\Domain M\to\Domain M$ fonksiyonu;
iki-li $\Assign{\Obj+}M$ ve $\Assign{\Obj\times}M$ fonksiyonları
(ikisi de $\Domain M^2\to\Domain M$); ve iki-li
$\Assign{\Obj<}M\subseteq\Domain M^2$ bağıntısı.

Böyle bir yapının apaçık bir örneği şudur:
\begin{enumerate}
\item $\Domain N=\Nat$;
\item $\Assign{\Obj0}N=0$;
\item bütün $n\in\Nat$ için $\Assign{\Obj\prime}N(n)=n+1$;
\item bütün $n,m\in\Nat$ için $\Assign{\Obj+}N(n,m)=n+m$;
\item bütün $n,m\in\Nat$ için $\Assign{\Obj\times}N(n,m)=n\cdot m$;
\item $\Assign{\Obj<}N=\Setabs{\tuple{n,m}}{n\in\Nat,m\in\Nat,n<m}$.
\end{enumerate}
Bu biçimde tanımlanan~$\Lang L_A$ yapısı~$\Struct N$,
$\Lang L_A$ dilinin mantıksal olmayan sabitlerini tam beklendiği gibi
yorumladığından \emph{aritmetiğin standart modeli} olarak adlandırılır.

Fakat~$\Lang L_A$ için başka birçok olası !!{structure}s vardır.
Örneğin evren olarak~$\Nat$ yerine tam sayılar kümesi~$\Int$'i alıp
$\Obj0$, $\Obj\prime$, $\Obj+$, $\Obj\times$ ve $\Obj<$ simgelerinin
yorumlarını buna göre tanımlayabiliriz. Bunun yanında~$\Lang L_A$ için
sayılarla uzaktan yakından ilgisi olmayan yapılar da tanımlayabiliriz.
\end{ex}

\begin{ex}
Küme kuramının~$\Lang L_Z$ dili için bir~$\Struct M$ yapısı yalnızca
bir küme ve tek bir iki-li bağıntı gerektirir. Dolayısıyla örneğin
insanlar kümesiyle ``$x$,~$y$'den yaşlıdır'' bağıntısı ya da~$\Nat$
ile $n,m\in\Nat$ için $n\geq m$ bağıntısı,~$\Lang L_Z$ için
!!a{structure} olarak kullanılabilir.

$\Lang L_Z$ için özellikle ilginç bir !!{structure}, evrendeki
!!{element}s öğelerinin gerçekten kümeler olduğu ve~$\Obj\in$
yorumunun gerçekten ``$x$,~$y$'nin !!a{element} öğesidir'' bağıntısı
olduğu \emph{kalıtsal sonlu kümeler} !!{structure}~$\Struct{{HF}}$
yapısıdır:
\begin{enumerate}
\item $\Domain{{HF}}=\emptyset\cup\Pow\emptyset\cup
  \Pow{\Pow\emptyset}\cup\Pow{\Pow{\Pow\emptyset}}\cup\dots$;
\item $\Assign{\Obj\in}{{HF}}=
  \Setabs{\tuple{x,y}}{x,y\in\Domain{{HF}},x\in y}$.
\end{enumerate}
\end{ex}

\begin{digress}
Neyin !!a{structure} sayıldığına ilişkin kabullerimiz mantığımızı
etkiler. Örneğin boş evrenleri dışlama seçimi, niceleyicili tümceler
için olağan sağlama (ya da doğruluk) açıklaması verildiğinde
$\lexists[x][(!A(x)\lor\lnot !A(x))]$ ifadesinin geçerli, yani
mantıksal bir doğru olmasını güvence altına alır. Bütün !!{constant}s
ifadelerinin evrendeki bir nesneye gönderme yapmasını istemek de
varlıksal genellemenin sağlam bir çıkarım kalıbı olmasını sağlar:
$!A(a)$, dolayısıyla $\lexists[x][!A(x)]$. Adların evren dışındaki
nesnelere gönderme yapmasına ya da hiç gönderme yapmamasına izin
verseydik, varlıksal genellemenin ek bir öncül gerektirdiği
\emph{serbest mantığa} doğru ilerlerdik: $!A(a)$ ve
$\lexists[x][\eq[x][a]]$, dolayısıyla $\lexists[x][!A(x)]$.
\end{digress}

